\scanpage{21}\mspage{20}
vrai pour tt $b$, alors $\unsure{L|\mathscr{L}}$ est \unsure{ttt|tt} ordonné. Plus \unsure{précisém$^{\mathrm t}$|précisément|généralem$^{\mathrm t}$}, pour que $L$ ttt ordonné, il f et s. que pour tt $b\in L$, $b$ soit comparable à tt autre élément, ce qui s'écrit
\[
\underbrace{L_{\leq b}\underset{\mathrm{mod}}{\cup}L_{\geq b}}_{L}=L_{\leq b}\cup L_{\geq b}
\]
\begin{center}\rule{2.5cm}{0.35pt}\end{center}

Partie loc. fermée modérée

\uline{Th} Soi\deleted{\unsure{ent}}t $X\deleted{\unsure{\subset|\supset}},Y$ deux parties fermées modérées, avec $X\supset Y$ (donc $X$ et $Y$ en pos. relative modérée)
\deleted{Considérons les} \deleted{Alors $\exists$ une plus petite partie fermée modérée \inserted{$Z$} contenant $X\setminus Y$.}
Soit
\[
\deleted{\unsure{\complement_X|C_X|\partial X}} C_{\mathrm{mod}}(Y,X)=\left\{x\in X\ \middle|\ 
\begin{aligned}&x\notin Y,\ x\text{ en pos. rel.}\\[-2pt]&\hphantom{x\notin Y,\ x\text{ en }}\text{mod. avec }Y\end{aligned}\right\}
\]
Alors $\exists$ plus petit ens. fermé modéré $Z$ \deleted{tel} \inserted{\unsure{contenant}} \deleted{\unsure{tt|te}} $C_{\mathrm{mod}}(Y,X)$, \deleted{\unsure{a}}\inserted{Ce} $Z$ \inserted{est $\subset X$, et} contient \deleted{\unsure{tt}} partie fermée modérée $T$ de $L$ telle que $T\subset X$, $T\cap Y=\emptyset$, $T$ en pos. rel. mod par. \inserted{\unsure{rapt|rel.}} à $Y$, \inserted{\unsure{De plus}} $Z$ est en position modérée avec $Y$.
% slanted
\margin{\uline{On pose} $Z=X\underset{\mathrm{mod}}{\setminus}Y$}
